\SourceLeaf{159}
Setzt man also $\alpha A^2+\beta B^2+\gamma C^2+\cdots$ als Summe an, so vernachlässigt man alle $(a-A)^2+(a'-A)^2+\cdots$, $(b-B)^2+(b'-B)^2+\cdots$, $\cdots$, die sämtlich positiv sind.

Der Beweis ist einfach. Es gilt nämlich
\[
\begin{aligned}
aa &= [A+(a-A)]^2 = AA+2A(a-A)+(a-A)^2\\
a'a' &= [A+(a'-A)]^2 = AA+2A(a'-A)+(a'-A)^2\\
a''a'' &= [A+(a''-A)]^2 = AA+2A(a''-A)+(a''-A)^2
\end{aligned}
\]
\begin{center}
. . . . . . . . . . . . . . . . .
\end{center}
\[
\Sigma aa = \alpha AA+2A\Sigma(a-A)+\Sigma(a-A)^2
\]
oder, da $\Sigma(a-A)=\Sigma a-\alpha A=0,$
\[
\Sigma aa = \alpha AA+\Sigma(a-A)^2.
\]
\Needspace{6\baselineskip}
\begin{center}
\textsc{Gauss} an \textsc{Olbers}. Göttingen, 3. Mai 1827.
\end{center}

\begin{center}
. . . . . . . . . . . . . . . . .
\end{center}

Für eine erfolgreiche Anwendung der Wahrscheinlichkeitsrechnung auf Beobachtungen ist umfassende \textls{Sach}kenntnis stets von höchster Wichtigkeit. Wo sie fehlt, ist es immer bedenklich, Beobachtungen wegen größerer Abweichungen auszuschließen, sofern ihre Zahl nicht sehr groß ist. Alle einzelnen Bestandteile des Beobachtungsfehlers, deren Vermeidung außerhalb unserer Möglichkeiten liegt, haben gewiss Grenzen, auch wenn wir sie nicht genau angeben können. In sehr vielen Fällen können wir mit Gewissheit sagen, dass ein aufgetretener großer Fehler außerhalb der Grenzen der Möglichkeit solcher Fehler liegt und ein außergewöhnliches Versehen begangen worden sein muss. Solche Fälle muss man natürlich ausschließen. Solange man sich aber vorstellen kann, dass der Fehler durch ein unglückliches Zusammentreffen der Bestandteile entstanden ist, soll man ihn nicht ausschließen. Gelegentlich kann es auch Fälle geben, bei denen zweifelhaft ist, ob sie zur ersten oder zur zweiten Klasse gehören. Dann mag man verfahren, wie man will,
